% ID: SEL-D0022-OA-E03
% ESTADO: cotejada. Segunda pasada visual completa y borrador compilado revisados.
% Fuente: Selectividad/matematicas_aplicadas/reserva_a_EXAMEN 5 -ANDALUCÍA.pdf
% SHA-256: 668689642ca997cf817c9f9f6608d2835efe898e3162b050d7b22b51f225d363
% Página: 1; opción A; ejercicio original 3.

El 60\% de los jóvenes de una ciudad usa Facebook, el 80\% usa WhatsApp y el 4\% usa Facebook pero no WhatsApp.
\begin{enumerate}[label=\alph*)]
\item \textbf{(0.5 puntos)} Halle el porcentaje de jóvenes de esa ciudad que usa ambas aplicaciones.
\item \textbf{(0.75 puntos)} Calcule el porcentaje de esos jóvenes que usa WhatsApp pero no Facebook.
\item \textbf{(0.75 puntos)} Entre los jóvenes que usan WhatsApp, \textquestiondown{}qué porcentaje usa también Facebook?
\item \textbf{(0.5 puntos)} Los sucesos ``usar Facebook'' y ``usar WhatsApp'', \textquestiondown{}son independientes?
\end{enumerate}
